\chapter{Family--defect relative curvature and auxiliary-sign gate}

Предыдущая глава оставила две внешне близкие возможности получить
\(\tau_3(XX^T-|\Phi|^2I_3)^2\): relative/mapping-cone curvature или
auxiliary \(D\)-sector. Здесь они проверяются отдельно внутри уже
построенной real KO6 geometry. Результат уточняет terminology: ordinary
mapping cone и строгий \(SO(3)\)-adjoint \(D\)-term не дают требуемый
functional. Формальная Legendre polarization на self-adjoint части middle
matrix summand существует, но не является положительным classical
auxiliary elimination.

\section{Oriented differential и middle Hodge commutator}

Пусть oriented differential трёхузловой цепи равен
\begin{equation}
 d=
 \begin{pmatrix}
 0&0&0\\
 X&0&0\\
 0&\Phi I_3&0
 \end{pmatrix}.
 \label{eq:family-oriented-quiver-differential}
\end{equation}
Тогда middle block Hodge commutator имеет правильный знак:
\begin{equation}
 \mu_G=[d,d^\dagger]\big|_{V_G}
 =XX^T-|\Phi|^2I_3.
 \label{eq:family-middle-hodge-commutator}
\end{equation}
Это identity oriented complex, а не identity ordinary self-adjoint
spectral quartic.

\section{Почему mapping-cone route не работает}

Возьмём canonical graph coordinate
\begin{equation}
 h=\operatorname{diag}(-I_3,0_3,I_3)
\end{equation}
и even curvature
\(F=(d+d^\dagger)^2\). Как и в Pati--Salam relative gate, functional
\begin{equation}
 S_{\rm mc}=\left\|\frac12[h,F]\right\|_{\rm HS}^2
 \label{eq:family-mapping-cone-functional}
\end{equation}
удаляет diagonal backtracking blocks и сохраняет endpoint block. Но теперь
endpoint composition равна \(\Phi X\), поэтому
\begin{equation}
 S_{\rm mc}=2|\Phi|^2\Tr XX^T.
 \label{eq:family-mapping-cone-endpoint-product}
\end{equation}
На radial slice это \(6\rho^2r^2\): присутствует только положительный
mixed monomial, а \(\rho^4\) и \(r^4\) отсутствуют.

\begin{nogocriterion}[Family mapping-cone no-go]
Canonical relative derivation, применённая к ordinary even curvature
\((d+d^\dagger)^2\), выбирает length-two endpoint composition, а не
middle commutator~\eqref{eq:family-middle-hodge-commutator}. Поэтому
Pati--Salam mapping-cone mechanism нельзя перенести в family branch простой
заменой block dimensions.
\end{nogocriterion}

\section{Почему это не строгий SO(3) D-term}

Для real locking field \(X\) и scalar arrow \(\Phi I_3\) matrix
\(\mu_G\) симметрична. Но gauge Lie algebra middle node равна
\begin{equation}
 \mathfrak{so}(3)=\operatorname{Skew}_3(\mathbb R).
\end{equation}
Ортогональная projection требуемой curvature на adjoint sector тождественно
исчезает:
\begin{equation}
 \operatorname{pr}_{\mathfrak{so}(3)}\mu_G
 =\frac12(\mu_G-\mu_G^T)=0.
 \label{eq:family-so3-dterm-zero}
\end{equation}
Следовательно, standard supersymmetric terminology ``\(SO(3)\) D-term''
здесь была бы некорректна. В обычной quiver moment-map literature node
groups являются unitary/Kähler groups, и real moment map принимает значения
в dual их Lie algebra; см. \path{arXiv:hep-th/0405230} и
\path{arXiv:0807.4734}. Наш real orthogonal reduction меняет этот вывод.

\begin{nogocriterion}[Strict adjoint D-term no-go]
Нельзя объявлять \(\tau_3(\mu_G^2)\) результатом elimination обычного
\(SO(3)\)-adjoint auxiliary field: adjoint projection самой \(\mu_G\)
равна нулю. Переход к \(U(3)\) исправил бы representation mismatch, но
нарушил бы уже установленный запрет на новую family unitary gauge algebra.
\end{nogocriterion}

\section{Self-adjoint polarization и sign obstruction}

У real matrix summand имеется каноническое self-adjoint пространство
\begin{equation}
 \operatorname{Sym}_3(\mathbb R)
 =\{K\in M_3(\mathbb R):K^T=K\}.
\end{equation}
Оно не является gauge Lie algebra, но сохраняется при
\(SO(3)\)-conjugation. Более того, linearization
\begin{equation}
 \delta\mu_G
 =\delta X+\delta X^T-2\,\delta r\,I_3
\end{equation}
в точке \(X=I_3\), \(|\Phi|=r=1\) имеет rank шесть и заполняет всё
\(\operatorname{Sym}_3(\mathbb R)\).

Hilbert--Schmidt polarization даёт точную algebraic identity
\begin{equation}
 \boxed{
 \tau_3(\mu_G^2)
 =\sup_{K=K^T}
 \left\{2\tau_3(K\mu_G)-\tau_3(K^2)\right\}.}
 \label{eq:family-self-adjoint-auxiliary-legendre}
\end{equation}
Stationary equation единственна:
\begin{equation}
 K_\star=\mu_G.
\end{equation}
Действительно, для любого trial \(K\)
\begin{equation}
 \tau_3(\mu_G^2)
 -2\tau_3(K\mu_G)+\tau_3(K^2)
 =\tau_3((K-\mu_G)^2)\ge0.
 \label{eq:family-auxiliary-legendre-gap}
\end{equation}
Однако functional в фигурных скобках имеет Hessian
\begin{equation}
 -2I_6,
\end{equation}
то есть stationary point является максимумом. Положительный real
Euclidean auxiliary action имеет противоположный результат:
\begin{equation}
 \inf_{K=K^T}
 \left\{\tau_3(K^2)-2\tau_3(K\mu_G)\right\}
 =-\tau_3(\mu_G^2).
 \label{eq:family-real-auxiliary-wrong-sign}
\end{equation}
Следовательно, formal supremum нельзя интерпретировать как elimination
устойчивого real classical field.

Положительный quartic допускает точное Gaussian representation только с
imaginary coupling:
\begin{equation}
 \tau_3(K^2+2iK\mu_G)
 =\tau_3((K+i\mu_G)^2)+\tau_3(\mu_G^2).
 \label{eq:family-imaginary-hs-completion}
\end{equation}
Интеграционный contour может оставаться real, но integrand становится
complex, а saddle лежит при \(K=-i\mu_G\). Это стандартная
Hubbard--Stratonovich развилка для repulsive channel; см.
\path{arXiv:2308.05150}. Поэтому данный механизм относится к quantum
measure, а не к уже выведенному classical parent action.

\section{Единая trace normalization}

Тот же normalized middle trace задаёт kinetic terms:
\begin{equation}
 S_{\rm kin}
 =\tau_3(DX(DX)^T)
 +\tau_3((D(\Phi I_3))^\dagger D(\Phi I_3))
 =\tau_3(DX(DX)^T)+|D\Phi|^2.
 \label{eq:family-common-kinetic-trace}
\end{equation}
В identities~\eqref{eq:family-self-adjoint-auxiliary-legendre} и
\eqref{eq:family-imaginary-hs-completion} используется ровно та же
\(\tau_3\). Отдельный continuous coefficient не возникает. Но symmetry не
фиксирует metric полностью: как \(SO(3)\)-module
\begin{equation}
 \operatorname{Sym}_3(\mathbb R)
 =\mathbb RI_3\oplus\operatorname{Sym}_{3,0}(\mathbb R)
 \simeq \mathbf1\oplus\mathbf5,
\end{equation}
поэтому invariant positive metrics имеют два независимых веса. Именно
полный represented calculus, а не одна \(SO(3)\)-covariance, должен был бы
вывести их равенство. Все algebraic identities и двухмерность metric
commutant проверены машинно.

\begin{theorem}[Auxiliary polarization without classical pass]
На middle summand существует coefficient-free polarization identity на
\(\operatorname{Sym}_3(\mathbb R)\), но её real stationary point является
максимумом. Устойчивый real classical auxiliary даёт противоположный знак;
правильный знак требует imaginary Hubbard--Stratonovich coupling и потому
остаётся quantum-measure construction.
\end{theorem}

\section{Точный оставшийся разрыв}

Необходимо проверить, возникает ли
\(\operatorname{Sym}_3(\mathbb R)\) как canonical self-adjoint represented
degree-two curvature module после quotient по junk. Если нет, imaginary
HS representation должна быть выведена из determinant/Pfaffian measure с
заранее фиксированным contour; простое объявление complex auxiliary не
считается parent derivation.

\begin{protocol}[Degree-two and measure continuation]
Следующий gate должен вычислить represented
\(\Omega_D^2(\mathcal A_{\rm q})\), его junk ideal и middle self-adjoint
quotient. При отрицательном результате остаётся проверить только imaginary
HS contour как следствие полной fermionic measure, а не как независимую
аксиому.
\end{protocol}

Итак, три classical реализации закрыты или условны: mapping cone выбирает
не тот block, strict \(SO(3)\) moment map равен нулю, а real symmetric
auxiliary имеет неверный sign. Surviving route локализован до represented
degree-two quotient либо complex quantum measure.

Следующий explicit junk audit закрывает первую возможность: на generic
фонe particle middle \(M_3\) целиком лежит в degree-two junk, а полный KO6
quotient сохраняет только одну central complex direction на conjugate
chain. Ни \(\operatorname{Sym}_{3,0}\), ни полный six-dimensional
\(\operatorname{Sym}_3\) module не выживают. Поэтому после следующей главы
условно открытым остаётся только imaginary HS measure или новая
differential-calculus architecture.

Машинный ledger сохранён в
\path{s2t_v4_family_defect_relative_auxiliary_moment_gate_results.json};
генератор ---
\path{s2t_v4_family_defect_relative_auxiliary_moment_gate.py}.